\documentclass[twocolumn]{article}
\usepackage[margin=1in]{geometry}
\usepackage[english]{babel}
\usepackage[T1]{fontenc}
\usepackage{cuted}
\usepackage[hyphens]{url}
\usepackage{hyperref}
\usepackage{graphicx}
\usepackage{ifthen}
\usepackage{multicol}
\usepackage{listings}
\usepackage[
    backend=biber,
    style=alphabetic,
]{biblatex}
\usepackage{enumitem}
\usepackage{appendix}
\usepackage{titling}
\usepackage[labelfont=bf]{caption}
\usepackage{setspace}
\usepackage{draftwatermark}
\usepackage{todonotes}
\usepackage{amssymb}
\addbibresource{main.bib}
\SetWatermarkText{Draft}
\onehalfspacing
\interfootnotelinepenalty=10000

\newcommand{\stefan}{Stefán Ólafsson}
\newcommand{\jacky}{Jacky Mallet}
\newcommand{\rem}{T-701-REM4}
\newcommand{\secu}{T-742-CSDA}

\newcommand{\titleHeading}{CryptoChronicles}
\newcommand{\titleSubheading}{LessPM's Quest to a Passwordless Utopian Ecosystem}
\newcommand{\subjectName}{\secu}
\newcommand{\instructorName}{\jacky}

\title{\titleHeading}
\author{Håvard Nordlie Mathisen}
\date{January, 2023}

\iffalse

"Breaking the Chains: Unleashing LessPM's Passwordless Revolution"
"Pass-go Paradox: The Magical Mystery of LessPM's Passwordless Power"
"No More Hide-and-Seek: Discovering LessPM's Cryptographic Wonderland"
"Keyless Kingdom: A Journey into LessPM's Ultra-Secure Realm"
"The Fort Knox of Password Managers: Unlocking LessPM's Passwordless Potential"
"Less is More: The Art of Mastering Passwordless Security with LessPM"
"CryptoChronicles: The Adventures of LessPM's Passwordless Pioneers"

"Phantom Passcodes: LessPM's Invisible Path to Ironclad Security"
"Unraveling the Enigma: Exploring LessPM's Passwordless Universe"
"Now You See Me, Now You Don't: LessPM's Passwordless Sleight of Hand"
"Cryptic Conundrums: Solving the Passwordless Puzzle with LessPM"
"The Digital Decoder Ring: Unlocking the Secrets of LessPM's Passwordless Mastery"
"Mission: Impeccable – LessPM's Quest for Ultimate Passwordless Security"
"Pass-Free Paradise: The Utopian World of LessPM's Passwordless Ecosystem"

\fi

\newboolean{draftmode}
\setboolean{draftmode}{true}

\setlength{\columnsep}{20pt}
\setlength{\skip\footins}{20pt}
\begin{document}
    \begin{titlepage}
        \begin{center}
            \includegraphics[scale=0.4]{images/HR_logo_hringur_transparent}
            % \vspace*{1cm}

            \huge
            \textbf{\titleHeading}\\
            \Large
            \textbf{\titleSubheading}

            \vspace{1cm}
            \large
            \subjectName

            \vspace{0.5cm}
            \Large
            Håvard Nordlie Mathisen\\
            \large
            Reykjavik University\\
            havard22@ru.is

            % \vspace{1.5cm}

            \vfill

            \normalsize
            \textit{Instructor}\\
            \instructorName

            \vspace{0.8cm}

            Department of Computer Science\\
            Reykjavik University\\
            March, 2023\\

        \end{center}
    \end{titlepage}


    \iffalse

    Purpose:
    Executive Summary: The primary goal of an executive summary is to present a condensed version of the report's main
    findings, conclusions, and recommendations. It is targeted at decision-makers or executives who may not have the
    time to read the entire report. The executive summary aims to persuade the reader to take action based on the
    presented information.

    Content:
    Executive Summary: The executive summary typically includes background information, the problem or issue addressed,
    methodology, key findings, conclusions, and recommendations.
    It may also emphasize the report's significance or the implications of the findings.

    LEngth:
    Executive Summary: Executive summaries are usually longer than abstracts, ranging from a few paragraphs to a
    few pages, depending on the length and complexity of the report.
    The executive summary aims to provide enough detail to give a comprehensive understanding of the report's key
    points while still being concise.

    \fi


    \ifthenelse{\boolean{draftmode}}{
        \begin{strip}
            \tableofcontents
        \end{strip}
    }


    \iffalse

    Abstract: The purpose of an abstract is to provide a brief, comprehensive
    overview of the report's content, including its objectives, methodology,
    results, and conclusions.
    It is primarily intended for researchers, academics, or other professionals who want to quickly assess whether the
    report is relevant to their interests before deciding to read the entire document.

    Abstract: An abstract usually contains a brief introduction to the topic, the research question or problem,
    the methodology used, a summary of the results, and the main conclusions drawn from the study.
    It does not typically include recommendations or opinions.

    Abstract: Abstracts are typically shorter than executive summaries, usually limited to a single paragraph or around
    150--300 words.
    The abstract provides a high-level overview of the report's content and is meant to be as concise as possible.

    \fi

    \section*{Abstract}\label{sec:abstract}

    Test\cite{lol}

    \addtocontents{TOC}{\protect\setcounter{tocdepth}{1}}
    \section{Introduction}\label{sec:introduction}
    Polybius' \textit{The Histories}~\cite{perseus_tufts} contains the first documented use of passwords, describing how
the Romans employed ``\textit{watchwords}'' to verify identities within the military.
This provided a transparent, simple way to allow or deny entry to restricted areas of authorized personnel only.
The story of secret writing (in this context referenced as cryptography) goes back the past 3000 years~\cite{history_cryptography_cryptanalysis},
where the need to protect and preserve privacy between two or more individuals blossomed.

Fernando J. Corbató is widely credited as the all-father of the first computer password when he was responsible for
the Compatible Time-Sharing System (CTSS) in 1961 at MIT~\cite{levy1984hackers}.
The system had a \texttt{"LOGIN"} command, which, when the user followed it by typing \texttt{"PASSWORD"},
had its printing mechanism turned off to offer the applicant privacy while typing the password~\cite{ctss_programmers_guide}.
Given the long history of passwords and their importance, one could argue that it was a natural and judicious
step in the evolution of computer systems.

In today's digital landscape, utilizing various identifiers (such as usernames, email addresses, or phone numbers)
combined with passwords has become a prevalent method for verifying an individual's identity and ensuring
their authorization to access restricted materials.

In 2004, a study titled \textit{"The Memorability and Security of Passwords"}~\cite{yan2000password} was conducted into
advising users on the entropy of passwords and ways someone can use to remember a or multiple passwords.
A typical standard for larger organizations with a form of password creation system is to emphasize the diversity of
smaller characters, capitalized characters, length, and not be commonly referred to in a dictionary~\cite{yan2000password}.
The study analyzed the effectiveness of different password-creation strategies, suggesting that acronym-based
passwords offer a delicate balance between memorability and security\cite{yan2000password}.

However, as technology has advanced, the limitations of password-based authentication have become increasingly
apparent, leading to the development of more sophisticated methods like Universal Authentication Framework
(UAF)~\cite{fido_uaf_overview} and WebAuthn\cite{webauthn_level_2} through the Fast IDentity Online Alliance
\href{https://fidoalliance.org}{(FIDO)} and The \href{https://www.w3.org}{World Wide Web Consortium} (W3C).

WebAuthn, short for Web Authentication, is an open standard for web-based authentication that enables users to
securely access online web services without relying on a traditional password.
Through a collaborative effort between FIDO and W3C, WebAuthn is developed to leverage asymmetric cryptography\footnote{
    Asymmetric cryptography uses a key pair consisting of a public key and a private key. The public key encrypts
    data, while the private key decrypts it. The keys are mathematically related, but deriving one from the other 
    is infeasible, ensuring secure communication and data exchange.}\footnote{
    \textbf{Note:} According to the library used to implement
    \href{https://docs.rs/jsonwebtoken/latest/jsonwebtoken/index.html}{jsonwebtokens in Rust}
    it is the private key that encrypts and the public key is responsible for \href{https://docs.rs/jsonwebtoken/latest/jsonwebtoken/struct.DecodingKey.html#method.from_rsa_pem}
    {decrypting}. \textit{Last Accessed: 2023-03-25}.}
and biometric or hardware-based authenticators to provide a more secure and robust authentication experience.

This report delves into the implementation of LessPM, a passwordless password manager that leverages WebAuthn to provide a secure
authentication experience, free from the constraints of traditional passwords, while placing a strong emphasis on
security.
By examining recent advancements in authentication mechanisms and the related innovative potential of WebAuthn, we
hope to illuminate the prospects of a passwordless future in digital security.


% \hyperref[sec:futurework]{security}

    \section{Background}\label{sec:background}
    
In today's world of rapidly evolving technology, it is essential to have a
strong foundation in the underlying concepts and protocols that drive modern
systems.
This background chapter aims to provide a comprehensive understanding of the key
technologies and principles that are relevant to our report.
By delving into these fundamental topics, we can better appreciate their
significance and application in the context of the implementation, design, and
architecture presented in the subsequent chapters.

We will provide background information on topics such as WebAuthn~\ref{subsec:webauthn},
JSON Web Tokens (JWT)~\ref{subsec:jsonwebtokens}, JSON Web Encryption (JWE)
\ref{subsec:jsonwebtokens}, hashing through Argon2~\ref{subsec:hashing-through-argon2},
and Advanced Encryption Standard (AES)~\ref{subsec:aes}.

% You should probably blabber some more here but we will do that when we have gotten further with the actual writing.

\subsection{WebAuthn}\label{subsec:webauthn}
We used WebAuthn to perform passwordless authentication in LessPM\@.
WebAuthn is a collaborative, open standard between the FIDO Alliance and W3C\@.
The aim of the standard is to implement a secure, robust key-based
authentication system for the web, to strongly authenticate users through
biometrics~\cite{webauthn_level_2}.

The concept relies on the use of a third-party device, called an
Authenticator Device (AD), which leverages asymmetric cryptography.
These devices employ biometric or hardware-based mechanism to provide secure
and reliable means of authenticating a user in LessPM\@.

Upon registering in LessPM, the AD generated a key-pair called a Passkey.
This Passkey contained a CID uniquely generated for each registered
key-pair~\cite{webauthn_credential_id,webauthn_public_key_credential}, per
services registered on the AD\@.

Further, if a malefactor gains access to an individual's Passkey, they might
compromise one specific service, whereas a traditional password could
potentially compromise multiple services where password reuse
occurs~\cite{wang2018next}.
This suggests that WebAuthn could create a stronger level of security, whereas
a traditional password exposed in a leak might leave a user susceptible to
losing access across services and devices through the reused password.
Preventing the potential security threat of password reuse is a top priority
for LessPM\@.
That is why WebAuthn has been chosen as the preferred standard for
passwordless user authentication in LessPM\@.

WebAuthn further attempts to mitigate phishing, man-in-the-middle, and
brute-force attacks through its intuitive design~\cite{webauthn_level_2}.
By leveraging the increasing market of smartphones with
biometrics~\cite{statista-biometric-transactions}, WebAuthn becomes a natural
extension, not just for LessPM, but for general authentication.

In order to increase the security of authenticated requests, each of the
steps using WebAuthn, Registration~\ref{subsubsec:metho-registration},
Authentication~\ref{subsubsec:metho-authentication} and
Password- Creation \& Retrieval~\ref{subsubsec:creation-and-retrieval})receives
different tokens with different information and is constructed, verified, and
discarded separately in LessPM.\footnote{
  The token for password creation/retrieval is the same as verifying that a user
  is signed in and authorized to create entries on behalf of the user.
}
We made this decision to further elevate security in LessPM, so that an
exposed token could not serve as more than one entry point for one of the steps.

\subsubsection{Registration}\label{subsubsec:metho-registration}
When a user tried registering in LessPM, the client performs a
registration request to the server, called a Relying Party (RP), carrying a
relevant User Identifier (UID. i.e.\ username, phone number, email, Etc.).
This is done to check whether a user with a similar UID existed\footnote{
  WebAuthn does not require to check whether the CID generated by the AD exists.
  However, accepting users with similar identifier might leave a risk of
  providing unauthorized access.
  Hence why LessPM performs such a check.
} in the system, and LessPM denies the registration process if it does.

If there are no users with a similar UID, LessPM responds by initiating a
Registration Ceremony (RC), generated a challenge and a Unique User
Identifier (UUID), which serves as the body of the HTTPS response.
LessPM also attached a JWE to the HTTPS response \texttt{Authorization} header
which is sent to the client.

The \texttt{expiration} time for this claim is set to one minute to allow the
user some time to authenticate.\footnote{
  WebAuthn describes a timeout performed within the system. In our case, this
  timeout is one minute.
  However, we chose to add this extra step to secure LessPM further, and
  JWEs require an expiry time (See Section~\ref{subsec:auth-and-auth}).
}
The claim expired after this minute, and the user would then have to restart the
process.
This expiration timer was a decision we made to emphasize security within LessPM
further.

Using the issued challenge, LessPM's client called the browser-integrated
WebAuthn API\footnote{
  In LessPM's case, this is the \textit{navigator.credentials} API provided by
  the browser.
}, prompting the user to utilize their AD to create a new Passkey credential
through the Client To Authenticator Protocol (CTAP2)\footnote{
  The user is prompted to use their AD to prove their presence, which can
  involve facial recognition, providing a fingerprint, or any other modality
  supported by the device that the user chooses.
} for
LessPM\@.
At this point, it is entirely up to the user to decide what device to use to
authenticate.
For our development, we used an iPhone 13 Pro Max, a Samsung S21, and a
Samsung Galaxy A52 to test authentication.\footnote{
  Other alternatives included a YubiKey, NitroKey, Etc.
}
We scanned the QR codes prompted through our phones, which initiated the
key-pair generation on the AD after a successful biometric scan, such as
facial recognition.\footnote{
  There is a question of concern that photography can bypass facial recognition.
  Apple uses built-in sensors to scan depth, colours and a dot projector to
  create a 3D scan of a person's face.
  However, this approach prevents the use of photography to authenticate through
  their FaceID and TrueDepth technology~\cite{apple-support}.
}

Finally, the AD signs the challenge using the private key stored on the
device.
The created public key, signed challenge, and additional metadata are combined
into a public key credential object, which is forwarded to the client
through CTAP2 and then sent to LessPM in a new request.

Before the new request reaches the RP/LessPM, an authentication middleware
checks and validates the JWE, denying the request with an \texttt{Unauthorized}
HTTPS status code if the request is invalid.\footnote{
  Validity in this context means not timed out, tampered with, or similar.
}
If the RP/LessPM can validate the authenticity of the signed challenge and
public key through WebAuthn, the user was stored in the database, along with the
UUID generated.

Thus completing the RC\@.
This process can be seen in Figure~\ref{fig:webauthn}


\subsubsection{Authentication}\label{subsubsec:metho-authentication}
When a user wishes to perform authentication (commonly referred to as
\texttt{logging in}), much of the same procedure occurs in LessPM.
The user issued an authentication request in the LessPM's client, carrying
the UID the user used to register.
The client included this information in the body.

Upon receiving an incoming request, the RP/LessPM checked the database for
the containing UID\@.
LessPM immediately rejected the request, should the user not exist in the
database.
If LessPM found an associated user with the UID in the database, the server
responds by initiating and issuing an Authentication Ceremony (AC),
collecting the public key associated with the user from the database and
generating a new challenge.
Upon validation\footnote{
  Validation in this context only means that the key stored in the database
  is valid.
} of the public key, LessPM creates a new JWE, which, like registration, also
receives an \texttt{expiry} time of one minute for the user to authenticate,
attached to the HTTPS \texttt{Authorization} header.

The challenge is issued and LessPM's client then calls the
browser-integrated WebAuthn API again, prompting the use of the original AD
to validate and sign the challenge through CTAP2.
Unlike the registration process, the AD now yields a signed signature based
on the challenge issued by the RP/LessPM\@, which is transferred back to
LessPM's client through CTAP2 with some AD specific
data~\cite{webauthn_authenticator_data}.
A new HTTPS request is then issued with the signed challenge and the AD specific
data, and the RP/LessPM then validated the signature using the stored public
key.
LessPM considers the user authenticated if the RP accepted and validated the
signed challenge.
If the AC is unsuccessful at any point in the ceremony, the ceremony is
aborted and considered invalid.

Upon success a new JWE is generated.
This time, the \texttt{expiry} time-to-live is set to a 15-minute\footnote{
  The specification does not specify any upper- or lower bounds for the
  expiry time~\cite{RFC7519}
}
timeframe, allowing the user some time to perform necessary activities
within LessPM\@.
The WebAuthn-related process can be seen in Figure~\ref{fig:webauthn} and the
JWE-related process can be seen in Figure~\ref{fig:JWT-process}.

\subsubsection{Password Creation \& Retrieval}\label{subsubsec:creation-and-retrieval}
Passwords are sensitive in nature, so it seems only natural in a security
context to enforce an extra level of authentication upon retrieving and creating
one unique password in LessPM.\footnote{
  We retrieved a complete list of the user's passwords upon successful authentication.
  The hashed version of the password is stripped of the returned values to protect and enforce security.
}

The following options are presented to a user when they initiate a password
creation process in LessPM's client:

\begin{itemize}
  \item \textbf{User Identifier}: An identification that the user wants to
  associate with the password entry they are storing.
  Such as a username, phone number, or email.
  \item \textbf{Website}: A URL or similar where the password belongs.
  \item \textbf{Password}: The user is prompted with the input to create a
  strong password automatically, choosing options such as numbers, special
  symbols, smaller or larger characters, and the length.
  As an option, the user was also permitted to construct their password but
  warned by a warning saying that this option is less secure.
\end{itemize}

As a final step before a password is created and stored, the user is prompted to
reauthenticate with their AD\@.
This is done in a similar manner as described for registration (See
Section~\ref{subsubsec:metho-registration}) and authentication (See
Section~\ref{subsubsec:metho-authentication}).

To retrieve the plaintext version of the password, LessPM enforces a new
authentication request through the AD\@.
We made this decision to attempt to ensure the owner\footnote{
  In this context, we distinguish between \texttt{user} and \texttt{owner}
  as the person that created the password, not the person currently using
  LessPM.
} of
the password's presence\footnote{
  This approach is also used in password managers on phones to avoid situations
  where a user might have left their computer unlocked.
}.
In such a scenario where the owner left their computer open while logged in to
LessPM, a malefactor would not be able to simply retrieve a password at will.
To decrypt a password, LessPM also required the CID from the AD to perform
the key-reconstruction for the encryption algorithm employed in LessPM\@.
Further details about how encryption is achieved can be seen in
Section~\ref{subsec:password-encryption}.

\subsection{Advanced Encryption Standard with 256-bit}\label{subsec:aes}
The Advanced Encryption Standard (AES) is a symmetric\footnote{
  Symmetric in this context refers to the key being the same to encrypt and
  decrypt.
} key encryption algorithm.
Since its inception in 1998, it has become the gold standard to encrypt various
information across applications~\cite{schneier2015applied,rijndael_book}, being adopted as the
successor of Data Encryption Standard (DES) by The National Institute of 
Standards and Technology (NIST) in 2001~\cite{nist_aes_winner}.
AES operates on a fixed-sizes units of data referred to as \texttt{blocks}~\cite{nistfips197blocks},
supporting keys of sizes 128-, 192-, and 256-bit~\cite{nistfips197intro}.
The design is based on a Substitution-Permutation Network (SPN) structure, which
combines substitution and permutation to provide a high level of security
through multiple rounds of processing~\cite{nistfips197specification}.
AES with 256-bit key-length (henceforth referred to as AES-256 in the rest
of the document), employs a 256-bit key and consists of 14 rounds of encryption,
offering an advanced level of security compared to its counterparts with shorter
key lengths and fewer rounds~\cite{nistfips197256}.
In each round of encryption, AES-256 undergoes four primary transformations:
\begin{figure}[htbp]
  \begin{itemize}
    \item \textbf{SubBytes} is a non-linear substitution step where each byte is replaced with
    another according to a predefined lookup table.
    \item \textbf{ShiftRows} is a transposition step where each row of the state
    is shifted cyclically a certain number of steps.
    \item \textbf{MixColumns} is a mixing operation that operates on the columns of the state,
    combining the four bytes in each column.
    \item \textbf{AddRoundedKey} combines the subkey with the state\protect\footnotemark
    ~using a bitwise exclusive OR (XOR) operation.
  \end{itemize}
  \caption{The steps AES takes when encrypting and decrypting data~\cite{nistfips197specification}.}
  \label{fig:aessteps}
\end{figure}
\footnotetext{The term \texttt{state} refers to an intermediate result that changes as
the algorithm progress through its phases.}
\newline
The larger key-size in AES-256 provides an exponential increase in number of
possible keys, making it significantly more resilient to brute-force attacks
and further solidifying its position as a robust encryption standard for
safe-guarding sensitive information.\footnote{
  The practical number of potential keys for an AES-256 implementation is
  $2^{256}$ possibilities. This gives us an approximation of $1.1579209 \times 10^{77}$
  options.
  The number is theoretical, as this is a worst-case scenario of options that an
  attacker has to go through in order to find the right key.
}



\subsection{Hashing through Argon2}\label{subsec:hashing-through-argon2}
Argon2 is a key-derivation function developed by Alex Biryukov, Daniel Dinu,
and Dmitry Khovratovich~\cite{argon2specs}.
Argon2 aims to provide a highly customizable function
tailored to the needs of distinct contexts~\cite{argon2specs}.
Additionally, the design offers resistance to both time-memory trade-off and
side-channel attacks as a memory-hard function~\cite{argon2specs}.

The function fills large memory blocks with pseudorandom data derived from the
input parameters, such as the password and salt.\footnote{
  A salt is a randomly generated sequence of characters, unique to each instance
  that gets hashed. Argon2's intension is to have a 128-bit salt for all
  applications but this can be sliced in half, if storage is a
  concern~\cite{argon2specs}.
}
The algorithm then processes the blocks non-linearly for a specified number of
iterations~\cite{argon2specs}.

The function offers three configurations, depending on the environment where the
function will run and what the risk and threat models are:

\begin{figure}[htbp]
  \centering
  \begin{itemize}
    \item \textbf{Argon2d} is a faster configuration and uses data-depending
    memory access.
    This makes it suitable for cryptocurrencies and applications with little to
    no threat of side-channel timing attacks.\protect\footnotemark
    \item \textbf{Argon2i} uses data-independent memory access.
    This configuration is more suitable for password hashing and key-derivation
    functions.\protect\footnotemark
    ~This configuration is slower due to making more passes over the memory as
    the hashing progresses.
    \item \textbf{Argon2id} is a combination of the two, beginning with
    data-dependent memory access before transitioning to data-independent
    memory access after progressing halfway through the process.
  \end{itemize}
  \caption{The three configurations of Argon2~\cite{argon2specs}.}
  \label{fig:argon2conf}
\end{figure}

\footnotetext{
  Side-channel timing attacks analyze execution time variations in cryptographic
  systems to reveal confidential data, exploiting differences in time caused by
  varying inputs, branching conditions, or memory access patterns.
}
\footnotetext{
  Due to the nature of prioritizing security, LessPM uses the third
  configuration. This is expanded upon in Section~\ref{subsec:hashing-aes-key}.}

Argon2, as a memory-intensive hashing function, demands substantial
computational resources from attackers attempting dictionary attacks.\footnote{
  A dictionary attack is an approach where an attacker tries to find a hash by
  searching through a dictionary of pre-computed hashes or generating hashes
  based on a dictionary commonly used by individuals or businesses.
}
This characteristic significantly hampers the feasibility of cracking passwords
using such attacks.
The algorithm's customizability allows users to adjust its behaviour based on
memory, parallelism, and iterations, catering to specific security requirements
and performance needs.
As these configurations are crucial for computing the original hash, Argon2
provides robust resilience against brute-force and side-channel attacks~\cite{
  argon2specs}.
The resulting enhanced security makes Argon2 suitable for password storage and
key-derivation in various applications and systems.

In 2015, Argon2 won the Password Hashing Competition~\cite{passwordhashing}.\footnote{
  NIST's competition to find an encryption algorithm inspired the Password
  Hashing Competition, but it took place without NIST's endorsement.
}


\subsection{JSON Web Token \& JSON Web Encryption}\label{subsec:json-web-tokens}\label{subsec:jsonwebtokens}
Although JSON Web Tokens (JWTs)~\cite{RFC7519} are not inherently encrypted,
they still serve an essential purpose for some forms of secure data transfer.
By combining JWTs with JSON Web Encryption (JWE)~\cite{rfc7516}, secure data
transfer between a client and a server can be achieved.

\subsubsection{JSON Web Token}
JSON Web Token (JWT) was introduced as part of RFC 7519 in 2015~\cite{RFC7519}.
It is a compact URL-safe string intended to transfer data between two entities.
They can be used as part of an authentication and authorization scheme in a web
service, application, or API\@.
The data in the string is intended as a payload and is referenced as a
\texttt{claim}~\cite{RFC7519}.

A JWT typically consist of three parts: header, payload, and signature.
The header and payload are serialized into JavaScript Object Notation (JSON)
and then encoded using a Base64Url encoding to ensure a URL-safe
format~\cite{RFC7519}.

The token contains an expiry timestamp, which, when decoded, is validated if
the timestamp is not passed at the time of decoding.
JWTs can be cryptographically signed using various algorithms like Hash-Based
Message Authentication (HMAC), Rivest-Shamir-Adleman (RSA), or
Elliptic-Curve Digital Signature Algorithm (ECDSA) ensuring the integrity and
authenticity of the token~\cite{RFC7519}.
This prevents unknown authorities from constructing or hijacking existing
tokens.\footnote{
  Hijacking a token could happen by a man-in-the-middle attack.
  This is done by a third-party individual listening and intercepting traffic
  in order to either read data or input their own in a client's request.
  This would allow an attacker to gain access to privileged information.
}

The URL-safe format of a JWT is often performed through a Base64 encoding, which
permits larger bits of data to be sent in a compressed, safe\footnote{
  Safe in this context should not be interchanged with secure.
  We reference safe as a way to transfer the data over the selected protocol,
  in most cases HTTP(S).
} format~\cite{RFC7519}.
JWTs are widely used for scenarios like single sign-on (SSO), user
authentication, and securing API endpoints by providing an efficient, stateless
mechanism for transmitting information about the user's identity, permissions,
and other relevant data~\cite{karande2018securingnode}.

The process functions as follows:
\begin{itemize}
  \item The client completed an authentication request.
  \item A token is constructed and created through a claim on the server.
  \item The claim can be any data the server wishes to use to authenticate the
  legitimacy of a future request.
  \item The claim gets signed with a desired algorithm.
  This can also be a secret, stored on the server.
  \item As part of the response to a request, the server appends the token.\footnote{
    A common place to embed these tokens is in the Authroization part of the
    HTTP header~\cite{RFC7519}.
  }
  \item The client receives the token and carries it upon the next performed
  request.
  \item When the server receives the token again, it validates the \texttt{exp}
  property of the JSON object and takes action accordingly.
\end{itemize}

\subsubsection{JSON Web Encryption}
Complimenting the JWT standard are JWEs.
There are many parallels between the two but the major distinction between them
are that JWEs are encrypted~\cite{rfc7516}.
This encryption can happen through AES (see Section~\ref{subsec:aes}) and
provides integrity and authenticity for the token.
This prevents eavesdropping or tampering with the token during transit
% This section might be split up into two parts, dividing JWT and JWE


    \section{Literature Review}\label{sec:literature-review}
    
The history of traditional passwords is a long and respectable one.
Polybius' \textit{The Histories}~\cite{perseus_tufts} contains the first
documented use of passwords, describing how the Romans employed
``\textit{watchwords}'' to verify identities within the military.
This provided a transparent, simple way to allow or deny entry to restricted
areas of authorized personnel only.
The story of secret writing (in this context referenced as cryptography) goes
back the past 3000 years~\cite{history_cryptography_cryptanalysis}, where the
need to protect and preserve privacy between two or more individuals blossomed.

Fernando J. Corbató is widely credited as the father of the first
computer password when he was responsible for the Compatible Time-Sharing
System (CTSS) in 1961 at MIT~\cite{levy1984hackers}.
The system had a \texttt{"LOGIN"} command, which, when the user followed it by
typing \texttt{"PASSWORD"}, had its printing mechanism turned off to offer
the applicant privacy while typing the password~\cite{ctss_programmers_guide}.
Given the long history of passwords and their importance, one could argue that
it was a natural and judicious step in the evolution of computer systems.

The study \textif{"The Memorability and Security of Passwords"}, conducted in
2004, provides insights into password creation strategies, including tips on
improving password entropy and methods for easy recall of passwords~\cite{
  yan2000password}.
With an emphasis on diversity in character selection, password length, and
avoiding dictionary words, the study suggests that acronym-based passwords offer
a delicate balance between memorability and security~\cite{yan2000password}.

However, as technology has advanced, the limitations of password-based
authentication have become increasingly apparent, leading to the development of
more sophisticated methods like Universal Authentication Framework
(UAF)~\cite{fido_uaf_overview} and WebAuthn~\cite{webauthn_level_2} through the
Fast IDentity Online Alliance \href{https://fidoalliance.org}{(FIDO)} and
The \href{https://www.w3.org}{World Wide Web Consortium} (W3C).

Passwordless Authentication (PA) is a growing field of study within Computer
Science, as traditional authentication methods like passwords are increasingly
recognized as vulnerable to attacks such as phishing\footnote{
  Phishing is a form of attack where a hacker tries to leverage Social
  Engineering to act as a trusted entity to dupe a victim to give away
  credentials by opening an email, instant message, or text message, then
  signing into a spoofed website, seeming legitimate~\cite{ripa2021emergence}.
} and credential stuffing.\footnote{
  Credential stuffing refers to the practice of using automated tools to
  inject compromised or stolen username and password combinations into web login
  forms to gain unauthorized access to user
  accounts~\cite{owasp-credential-stuffing}.
}
Passwords and sensitive information can also be a victim of successful
brute-force attacks~\cite{bonneau2012science} through data leakages by hacking
or purchasing information on the dark Web.

Following NIST~\cite{NIST:SP:800-171r2, NISTSP800-63-3}, authentication should
consist of covering one of the three principles in
Figure~\ref{fig:secprinciples}.

\begin{figure}[htbp]
  \begin{itemize}
    \item \textbf{Something you know}:
    Such as a password, an answer to a personal question, or a Personal
    Identification Number (PIN).
    \item \textbf{Something you have}:
    A device that contains some token or cryptographically signed keys.
    \item \textbf{Something you are}:
    Biometrics of any sort or kind.
    Facial recognition, retina scan, fingerprint and similar.
  \end{itemize}
  \caption{
    The Three Principles of Password Security~\cite{schneier2000secrets, NIST:SP:800-171r2}.
  }
  \label{fig:secprinciples}
\end{figure}

There are many approaches to passwordless authentication or a second
step to authenticate with a one-time password.\footnote{
  Often referenced as Two-Factor Authentication or Multi-Factor Authentication.
}
In 2022, Parmar et al.\cite{parmar2022} described several attractive solutions,
along with their advantages and drawbacks, for performing PA using the most
common methods.
The study discovered that PA commonly gets accepted as the most frictionless
authentication system for User Interfaces (UI)~\cite{parmar2022}.
Biometrics was mentioned as one of the authentication methods, concluding that
it captures universal human traits, encouraging differentiation from one
another~\cite{parmar2022}.
The same study raises the caution surrounding the loss of authentication
devices and how fingerprints can be compromised~\cite{parmar2022}.\footnote{
  The security implication of using the core concept of FIDO2's WebAuthn is
  subject to storage in the system on Apple-specific
  devices~\cite{appleSecureEnclave}.
  On an Android device, the implementation is up to the manufacturer of the
  device, where Samsung has implemented a Physically Unclonable
  Function~\cite{lee2021samsung} (PUF).
}

One promising approach is using the FIDO Alliance's collaborative work with W3C
to create WebAuthn.
WebAuthn permits users to authenticate through biometric information
stored on a device in the user's possession (i.e.\ phone, computer) or a
physical security key (i.e.\ YubiKey, Nitrokey, Etc.)~\cite{webauthn_level_2}.

In ~\cite{huseynov2022passwordless}, Huseynov utilized a Web interface with
WebAuthn to create credentials that users could use for a VPN\@.
The client required a user to authenticate through the procedure of WebAuthn.
On a successful request, the Remote Authentication Dial-In User Service (RADIUS)
creates a temporary username and password, which would then be transferred
as a response to the end-user, permitting them to copy and paste it into the
necessary client, alternatively to construct a batch file which would establish
the correct connection~\cite{huseynov2022passwordless}.
The study suggested creating a solution for a VPN client which embedded some
browser components~\cite{huseynov2022passwordless}.

Gordin et al.~\cite{gordin2021moving} implemented PA into an OpenStack
environment using WebAuthn, which increased security and bypasses the risk of
malefactors employing leaked passwords on other
services~\cite{gordin2021moving}.\footnote{
  See Section~\ref{subsec:webauthn-methodology} for further explanation as to how this
  works.
}
The PA process considerably reduced the number of attacks towards a server
because the server no longer has user authentication
secrets~\cite{gordin2021moving}.
They, however, utilized a fingerprint as the primary biometrics, citing cost as
a primary factor to discourage the use of
FIDO2~\cite{gordin2021moving}.\footnote{
  Authentication was provided through the Keystone environment on the OpenStack
  platform.
}

Statistia reported in 2020 that between 77--86\% of smartphones now have a form
of biometric scanner built into their
device~\cite{statista-biometric-transactions}.
Gorin, Et al.\ continue by mentioning that some individuals have trouble
using their biometric scanner or getting it to work correctly on their
device~\cite{gordin2021moving}, which can be a potential drawback for user
adaptability.

According to a study conducted by Lyastani et al.\ in
2020~\cite{ghrobany2020fido2}, a significant portion of users found the usage
of WebAuthn and Fido2 standard to be easy and secure, but with some concerns
about losing access to their accounts or fear of others accessing their
accounts~\cite{ghrobany2020fido2}.
The study utilized a fingerprint Yubikey for the authentication process.
Despite some reservations, the study found that overall, WebAuthn and Fido2
have good usability for passwordless authentication~\cite{ghrobany2020fido2}.

The authors reported that users automatically associated the loss of the AD with
losing access to the account~\cite{ghrobany2020fido2} -- and vice versa --
indicating that users are slightly unwilling to replace the initial principle of
\textit{"Something you know"} with the second \textit{"Something you have"} and
third \textit{"Something you are"} principle.
Additional research is necessary to educate users, increase trust and confidence
in the technology, and address concerns about the potential loss of account
access.\footnote{
  We believe that these concerns are mostly raised due to using a YubiKey
  and that using a phone-based authenticator would reveal other results.
}

Morii et al.~\cite{morii2017research} investigated the potential of FIDO as a
viable PA solution in 2017 when the FIDO2 and WebAuthn standards were yet to be
widely adopted.
The authors noted that, at that time, only the Edge browser had implemented
proper support for FIDO2 and WebAuthn~\cite{morii2017research}.
Despite the limited browser support, the study demonstrated the feasibility
of integrating PA into the well-established authentication system, Shibboleth
~\cite{shibboleth, morii2017research}.\footnote{
  Shibboleth is a widely-used, open-source federated identity solution that
  enables secure single sign-on across multiple applications and organizations.
}

As the technology evolved from FIDO to FIDO2.0, some security concerns
persisted, such as session hijacking\footnote{
  Session hijacking refers to an attacker gaining unauthorized access to a
  user's authenticated session, often exploiting weaknesses in the handling of
  cookies, sessions, or JSON Web Tokens (JWT).
  See Section~\ref{subsec:auth-and-auth}.
}, which can compromise
user accounts~\cite{morii2017research}, highlighting the need to protect these.

Since the publication of Morii et al.'s research, browser support for FIDO2 and
WebAuthn has significantly improved, with major browsers like Google Chrome,
Mozilla Firefox, Apple Safari, Microsoft Edge, and Vivaldi now offering
support for these standards.
This broader adoption has enabled the more widespread deployment of PA
solutions, providing increased security and improved user experiences across
various online services.

However, the ongoing evolution of security threats and the increasing
sophistication of attackers highlights the need for continuous research and
development in the field of passwordless authentication, ensuring that new
methods and standards are both secure and user-friendly.

Having explored various studies and developments in the field of PA, it is
evident that this area has been continuously evolving to provide secure,
user-friendly solutions.
However, the implementation of these solutions into stand-alone, real-world
applications, such as password managers, independent of other login systems, is
a critical aspect that requires thorough investigation.

In the following section, we will examine the methodology employed in this
report to integrate a passwordless system into a password manager,
taking into account the challenges and concerns identified in the literature.
By doing so, we aim to contribute to the growing body of knowledge on
passwordless authentication and its practical applications.


\iffalse
% More lighthearted reading.
Here are some additional studies and research papers related to passwordless authentication and related technologies. These studies can provide further insights and perspectives on the subject, expanding your knowledge and understanding of the field.

Biddle, R., Chiasson, S., & van Oorschot, P. C. (2012). Graphical Passwords: Learning from the First Twelve Years. ACM Computing Surveys, 44(4), 1-41. [DOI: 10.1145/2333112.2333114]

This survey paper provides a comprehensive review of graphical password schemes, an alternative authentication method that can be used in passwordless systems.
Bonneau, J., Herley, C., van Oorschot, P. C., & Stajano, F. (2015). Passwords and the Evolution of Imperfect Authentication. Communications of the ACM, 58(7), 78-87. [DOI: 10.1145/2699412]

This paper discusses the limitations of traditional password-based authentication and explores the potential of alternative authentication methods.
Gunson, N., Marshall, D., Morton, H., & Jack, M. (2011). User perceptions of security and usability of single-factor and two-factor authentication in automated telephone banking. Computers & Security, 30(4), 208-220. [DOI: 10.1016/j.cose.2011.01.006]

This study compares user perceptions of single-factor and two-factor authentication methods in the context of automated telephone banking systems, providing insights into the user experience of different authentication approaches.
Jakobsson, M., & Myers, S. (Eds.). (2007). Phishing and Countermeasures: Understanding the Increasing Problem of Electronic Identity Theft. John Wiley & Sons.

This book covers various aspects of phishing attacks and countermeasures, including passwordless authentication methods that can help mitigate the risks associated with phishing.
Karole, A., Saxena, N., & Christin, N. (2011). A Comparative Evaluation of Two-Factor Authentication Schemes. In Proceedings of the 28th Annual Computer Security Applications Conference (pp. 173-182). ACM. [DOI: 10.1145/2076732.2076759]

This paper presents a comparative evaluation of various two-factor authentication schemes, which can be useful for understanding the strengths and weaknesses of different passwordless authentication methods.
These additional studies can provide a broader understanding of the passwordless authentication landscape, offering further insights into alternative authentication methods, user perceptions, and the potential challenges and opportunities associated with implementing passwordless solutions.

\fi

    \section{Methodology}\label{sec:methodology}
    
Exploring the development and implementation of LessPM, our focus will be on the
key components, technologies, and steps that form the system's development
process.
A crucial part of the development and system was its security and robustness.

Our approach encompasses an explanation of the system architecture, the
technologies and tools utilized and the development process to create the
prototype.
By providing a comprehensive account of the system development process, we aim
to enable readers to understand the technical aspects of our work as well
as to assess the validity and relevance of our findings.

\subsection{Environment}\label{subsec:environment}
To implement the system, we approached the development from a type-safety
environment.

\subsubsection{Server}
We chose Rust as the programming language for our backend, which provided
significant benefits regarding the software development environment.

Rust's emphasis on safety and performance allowed us to create a highly
secure and efficient environment for our implementation without the risks
typically associated with memory-related
vulnerabilities~\cite{rivera2019preserving}.
The built-in memory management and focus on concurrency ensures that the
software runs smoothly, which is crucial~\cite{fischer1985impossibility} when
dealing with sensitive user data.

Additionally, Rust's surrounding ecosystem is at rapid
growth~\cite{librs-stats}, containing a vast library of high-quality
crates\footnote{
  \texttt{Crate} is the Rust-specific name for a package or library.
}, which enables expeditious development and easy integration of various
functionality.
Utilizing Rust's distinctive features, the LessPM achieves enhanced security and
reliability in the context of user authentication~\cite{rivera2019preserving},
showcasing the benefits of utilizing a modern programming language.

LessPM ran an instance of an HTTPS server, serving as a wrapper for the
sensitive data.\footnote{
  We took advantage of the Authorization header during development, as specified
  in RFC 7519.
  However, the framework we used for the server required us to expose the usage
  of \texttt{Authorization header} to access it in the client.
  See Section~\ref{subsec:security-analysis} for further explanation.
}
The Chromium developers mandate this constraint to guarantee that the pertinent
API is invoked exclusively within a secure context~\cite{webdev2021credential}.
During development, we established a secure connection by using a self-signed
certificate for the \texttt{``localhost''} domain.
As a deliberate decision, we opted for MongoDB due to its NoSQL architecture,
which facilitated the storage of Object-like structures~\cite{mongodb2021nosql}
including passwords, user accounts, and other related data that LessPM stores.

\subsubsection{Client}
An essential part of the implementation was to create a viable client that could
function as a visual entry point to LessPM and proof-of-concept.
For the simplicity of the project, we chose React as a framework.
React is a JavaScript library for building user interfaces, offering an
efficient and flexible approach to web development.
The design of LessPM was influenced by the principles of
least-knowledge~\cite{lieberherr1990assuring}.
This influenced the decision to ensure that the client only retained
necessary information to function properly, containing no knowledge of
the server nor its implementation.

We took advantage of React's ability to construct a single-page application with
no routing capabilities, avoiding the possibility of utilizing any URL tampering
or manipulation\footnote{
  Malefactors can perform URL Manipulation, which involves modifying a URL to
  request resources that would otherwise be inaccessible to a user.
} to attempt privilege escalation or
accessing restricted data.

When initiating the project's development, we had a strong vision of creating
a client that could seamlessly integrate with Chromium-based browsers as
an extension.\footnote{
  In the client's project folder are traces of a manifest.json file and build
  scripts to have the client run in an extension.
}
However, we quickly discovered that the Rust cargo used to perform WebAuthn
requests did not implement this functionality.
We reported this issue
\href{https://github.com/kanidm/webauthn-rs/issues/288}{upstream to the authors}
of the cargo, and we will continue to work with the authors to ensure a proper
implementation of this functionality in the future.

LessPM's client maintains authentication and authorization through an
encrypted JSON Web Token (JWT, RFC7519), inspired by JSON Web
Encryption (JWE, RFC7516), passed back and forth between the server and client.

\subsection{Authentication \& Authorization}\label{subsec:jwt}
A system that maintains authentication and authorization through a stateless
protocol such as HTTPS, requires some information to authorize a
client between requested resources.
LessPM took advantage of JWTs and an inspired variation of JWEs for this
purpose.
Although JWTs are not inherently encrypted~\cite{RFC7519}, combined with JWE,
they serve as an essential backbone for secure data transfer in LessPM\@.

\subsubsection{JSON Web Tokens}
JWT is a compact, URL-safe string intended to transfer data between two
entities, often represented as a Base64 encoding.
They are often used as part of authentication and authorization scheme within
a web service, application, or Application Programming Interface (API), and
LessPM does this as well.
The data in the string is delivered as a payload and is referred to as a
\texttt{claim}~\cite{RFC7519}.

The JWT typically consist of three parts: header, payload, and signature.
The header and payload are serialized into JavaScript Object Notation (JSON)
and then encoded using Base64 to ensure a URL-safe format.
This permits larger bits of data to be sent in a compressed, safe\footnote{
  Safe in this context should not be interchanged with secure.
  We reference safe as a way to transfer the data over the selected protocol,
  in LessPM's case this HTTPS.
} format.
JWTs are widely used for scenarios such as single sign-on (SSO), user
authentication, and securing API endpoints by providing an efficient and
stateless mechanism for transmitting information about the user's identity
and permissions~\cite{karande2018securingnode}.
According to~\cite{RFC7519}, these should be attached to the
\texttt{Authorization} part of the HTTPS header, along with a hardcoded
\texttt{Bearer} part, preceding the token as the value of the header.
\todo[inline]{
  Mention this in the security analysis.
}
In LessPM, the JWT claim contained data we could use to validate and
authenticate a user between HTTPS requests.

LessPM utilizes two functionalities of JWTs to ensure the validity of the
Base64-encoded token.

First, all tokens passed back and forth between the LessPM server and client
receives an expiry time.
We enforce this technique to ensure and maintain the security of accessed
resources, requiring a user to reauthenticate when the expiry timestamp is
passed.

Second, The tokens used in LessPM are cryptographically signed with RSA
Signature Scheme with Appendix - Probabilistic Signature Scheme (RSASSA-PSS)
using Secure Hashing Algorithms with 512-bit size (SHA-512), and is derived
from the Rivest-Shamir -Adleman (RSA) encryption.
The encryption is performed with a key-pair stored on the server to mitigate
the risk of token forgery and so LessPM can maintain the validity and
authenticity of the tokens created.
This prevents unknown authorities from constructing or hijacking tokens
belonging to LessPM\@.\footnote{
  Hijacking a token could happen by a man-in-the-middle attack.
  This is done by a third-party individual listening and intercepting traffic
  in order to either read data or input their own in a client's request.
  This would allow an attacker to gain access to privileged information.
}
Other signature algorithms that can be used include Hash-Based Message
Authentication (HMAC), Elliptic-Curve Digital Signature Algorithm (ECDSA), or
a plaintext secret stored on the server.
LessPM uses RSASSA-PSS with SHA-512 due to the asymmetric operations the
algorithms perform on the encryption and decryption, whereas HMAC and ECDSA is
based on symmetric- and elliptic curve cryptography.


\subsubsection{JSON Web Encryption}
By default, JWTs are not encrypted.
It is out of the question to unnecessarily expose any sensitive details about
the LessPM server's inner workings.
This lead to the decision of finding a way to encrypt the contents of the 
token before they are sent from LessPM's server.
Without the encryption, any attacker could decode the Base64 encoding and
gain access to the claim that LessPM uses to authorize and authenticate.

To perform this operation, LessPM seeks inspiration from JWEs and
RFC7516~\cite{rfc7516}.\footnote{
  To properly implement JWE, metadata should be encoded into the token, such as
  what algorithm performs the encryption~\cite{rfc7516}. For further
  explanation, see Section~\ref{sec:futurework}.
}
\todo[inline]{
  Diagram that shows the operation
}
Before the RSASSA-PSS signed JWT is attached to the HTTPS header,
LessPM encrypted the Base64-encoded data using AES-256, before turning them into
a new Base64 encoded string.
Encrypting the token with AES-256 elevated the overall security of the
JWTs, making them more suitable for transferring sensitive data between the
intercommunicating entities in LessPM\@.
As a final step, the ciphertext created by AES-256 is again encoded into a
new Base64.
When LessPM again receives this claim, we used it to gather details about the
user who made the request and responded to the request accordingly.
An example of this encoding can be seen in
Appendix~\ref{subsec:base64-and-ciphertext}.

\todo[inline]{
  Turn this into a diagram, because Jacky love those.
}
The process functions as follows in LessPM:
\begin{enumerate}
  \item The client completed an authentication request.
  \item A claim was constructed and created on the server
  The claim could have been any data the server wished to use to authenticate
  the legitimacy of a future request.
  \item The claim was signed with the desired algorithm.
  This could also have been a secret, stored on the server.
  \item The claim is encoded into a Base64 string.
  \item The Base64 encoded string is passed to AES-256 encryption algorithm,
  creating a ciphertext.
  \item The cipertext is created into a new Base64-encoding, which is the
  final token.
  \item As part of the response to a request, the server appended the
  token.
  \item The client received the token and carried it upon the next performed
  request.
  \item When the server received the token again, it decodes the
  Base64-encoded ciphertext.
  \item The ciphertext is decrypted, decoded and reconstructed to a claim.
  \item The \texttt{exp} property of the JWT is validated and LessPM takes
  action accordingly.
\end{enumerate}

\todo[inline]{
  Put this paragraph somewhere in the WebAuthn section
  Further, the server checked and verified the token before proceeding with any
  requests made to it.
  This information is updated and inspected between each
  re-render~\cite{react-component} of the client.
}

\todo[inline]{
This should go in the new section of WebAuthn in methodology. Somewhere.
For each step in registration~\ref{subsubsec:metho-registration},
authentication~\ref{subsubsec:metho-authentication}, and
password creation/retrieval~\ref{subsubsec:creation-and-retrieval}, different
tokens with different information are constructed, verified, and
discarded.\footnote{
  The token for password creation/retrieval is the same as verifying that a user
  is signed in and authorized to create entries on behalf of the user.
}
We opted for this approach so that an exposed token could not serve as more than
one entry point for one of the steps.
}



\subsection{WebAuthn}\label{subsec:webauthn-methodology}
We used WebAuthn to perform passwordless authentication in LessPM\@.
WebAuthn is a collaborative open standard between the FIDO Alliance and W3C\@.
The aim of the standard is to implement a secure, robust key-based
authentication system for the web, to strongly authenticate users through
biometrics~\cite{webauthn_level_2}.
The concept relies on the use of a third-party device, called an
Authenticator Device (AD), which leverages asymmetric cryptography.
These devices employ biometric or hardware-based mechanism to provide secure
and reliable means of authenticating a user in LessPM\@.

Upon registering in LessPM, the AD generates a key-pair called a Passkey.
This Passkey contains a CID uniquely generated for each registered
key-pair~\cite{webauthn_credential_id,webauthn_public_key_credential}, per
services registered on the AD\@.

Further, if a malefactor gains access to an individual's Passkey, they might
compromise one specific service, whereas a traditional password could
potentially compromise multiple services where password reuse
occurs~\cite{wang2018next}.
This suggests that WebAuthn could create a stronger level of security, whereas
a traditional password exposed in a leak might leave a user susceptible to
losing access across services and devices through the reused password.
Preventing the potential security threat of password reuse is a top priority
for LessPM\@.
That is why WebAuthn has been chosen as the preferred standard for
passwordless user authentication in LessPM\@.

WebAuthn further attempts to mitigate phishing, man-in-the-middle, and
brute-force attacks through its intuitive design~\cite{webauthn_level_2}.
By leveraging the increasing market of smartphones with
biometrics~\cite{statista-biometric-transactions}, WebAuthn becomes a natural
extension, not just for LessPM, but for general authentication.

\subsubsection{Registration}\label{subsubsec:metho-registration}
When a user tried registering in LessPM, the client performs a
registration request to the server, called a Relying Party (RP), carrying a
relevant User Identifier (UID. i.e.\ username, phone number, email, Etc.).
This is done to check whether a user with a similar UID existed\footnote{
  WebAuthn does not require to check whether the CID generated by the AD exists.
  However, accepting users with similar identifier might leave a risk of
  providing unauthorized access.
  Hence why LessPM performs such a check.
} in the system, and LessPM denies the registration process if it does.

If there are no users with a similar UID, LessPM responds by initiating a
Registration Ceremony (RC) and generated a challenge, a Unique User
Identifier (UUID) version 4, and a JWE which is sent to the client.

Using this challenge, LessPM's client called the browser-integrated WebAuthn
Application Programming Interface (API)\footnote{
  In LessPM's case, this is the \textit{navigator.credentials} API provided by
  the browser.
}, prompting the
user to utilize their AD to create a new Passkey credential
through the Client To Authenticator Protocol (CTAP2)\footnote{
  The user is prompted to use their AD to prove their presence, which can
  involve facial recognition, providing a fingerprint, or any other modality
  supported by the device that the user chooses.
} for
LessPM\@.

Finally, the AD signs the challenge using the private key stored on the
device.
The created public key, signed challenge, and additional metadata are combined
into a public key credential object, which is forwarded to the client
through CTAP2 and then sent to LessPM\@.

LessPM then verifies the authenticity of the signed challenge and public key
through WebAuthn, before storing the key with the generated UUID and the
username provided.
Thus completing the RC\@.
\todo[inline]{
  This process can be sign in figure huhu.
}

We generated a JWE during this process, which was signed with RSASSA-PSS using
SHA-512 before getting encrypted with AES-256
(See Section~\ref{sec:futurework}).
This claim is then attached to the request's \texttt{Authorization} header,
along with the Creation options from WebAuthn
(See Section~\ref{subsec:webauthn}).

The \texttt{expiration} time for this claim is set to one minute to allow the
user some time to authenticate.\footnote{
  WebAuthn describes a timeout performed within the system. In our case, this
  timeout is one minute.
  However, we chose to add this extra step to secure the project further, and JWEs require an expiry time. See Section~\ref{subsec:json-web-tokens}.
}
The claim expired after this minute, and the user would then have to restart the
process.
This expiration timer was a decision we made to emphasize security within LessPM
further.

When the response returned to the user, LessPM used the browser's built-in
WebAuthn API to prompt the user for their authentication
(See Section~\ref{subsec:webauthn}).
At this point, it is entirely up to the user to decide what device to use to
authenticate.
We used an Apple iPhone 13 Pro Max, a Samsung S21, and a Samsung Galaxy A52 to
test authentication.\footnote{
  Other alternatives included a YubiKey, NitroKey, Etc.
}

We scanned the QR codes prompted through our phones, which initiated the
key-pair generation on the device after a successful facial recognition.
\footnote{
  There is a question of concern that photography can bypass facial recognition.
  Apple uses built-in sensors to scan depth, colours and a dot projector to
  create a 3D scan of a person's face.
  However, this approach prevents the use of photography to authenticate through
  their FaceID and TrueDepth technology~\cite{apple-support}.
}
The public key then gets transmitted to the browser and sent to the Relying
Party through a new request.

Before the request reaches the Relying Party, an authentication middleware
checks for the JWE sent earlier and denies the request with an
\texttt{Unauthorized} HTTP status code if the request is invalid.
\footnote{
  Validity in this context means not timed out, tampered with, or similar.
}
If WebAuthn can validate the key and metadata sent with the request, the user
gets stored in the database, along with the UUID generated.

\subsubsection{Authentication}\label{subsubsec:metho-authentication}
Authentication worked similarly to the registration process.
Upon receiving an incoming request, the server checked the database.
The server immediately rejected the request if the user did not exist in the
database.

Further, the server attempted an Authentication Ceremony by collecting the
public key from the database.
Upon validation\footnote{
  Validation in this context only means that the key stored in the database is
  valid.
} of the public key, the server creates a new JWE, which also receives an
\texttt{expiry} of one minute for the user to authenticate (See
Section~\ref{subsec:webauthn}), before sending that and the challenge in the
response.
We then prompt the user to authenticate with their original authenticator
that they used in the registration step, having the authenticator sign the
challenge before performing a new request to the server with the signed
challenge.
The server considered the user authenticated if the Relying Party accepted the
signed challenge.

Finally, the user had a 15-minute timeframe to perform the necessary activity
within LessPM\@.
The specification for JWT does not specify upper- or lower bounds for this
value ~\cite{RFC7519}.

\subsubsection{Password Creation \& Retrieval}\label{subsubsec:creation-and-retrieval}
Passwords are sensitive in nature, so it seems only natural in a security
context to enforce an extra level of authentication upon retrieving and creating
one unique password.\footnote{
  We retrieved a complete list of the user's passwords upon successful authentication.
  The hashed version of the password is stripped of the returned values to protect and enforce security.
}

A user is presented with the following options when they initiate a password
creation:

\begin{itemize}
  \item \textbf{User Identifier}: An identification that the user wants to
  associate with the password entry they are storing.
  Such as a username, phone number, or email.
  \item \textbf{Website}: A URL or similar where the password belongs.
  \item \textbf{Password}: The user is prompted with the input to create a
  strong password automatically, choosing options such as numbers, special
  symbols, smaller or larger characters, and the length.
  As an option, the user was also permitted to construct their password but
  warned by a warning saying that this option is less secure.
\end{itemize}

As a final step before a password is created and stored, the user is prompted to
reauthenticate with their AD\@.
The password is created and stored in the database if this authentication
process is successful.

Users needed to make a new authentication request when they wanted to retrieve
the plaintext version of the password.
This step is done in part to validate the user's presence\footnote{
  This approach is also used in password managers on phones to avoid situations
  where a user might have left their computer unlocked.
} but also to gather the necessary bits for the key-reconstruction
(See Section~\ref{subsec:password-encryption}).

\subsection{Cors}\label{subsec:cors}
Cross-Origin Resource Sharing (CORS) must be configured correctly when the
server and client are running separately on different ports.

When a web page tries to access a resource hosted on another domain, browsers
perform an additional request to the server, called a \texttt{``preflight''}.
The preflight request determines whether the
request that the web page is trying to make to the server is allowed.
This request is done through the \texttt{OPTIONS} method in HTTP and contains
some information about the origin, accepted Content-Type, and similar for the
actual request.

The server responds to this with what methods and headers are allowed, denying
the actual request from ever happening if the preflight is not successful.

We constructed a CORS \texttt{layer}\footnote{
  In this context, we referred to a layer as a wrapper around all other
  requests.
} which contained the two domains for the server and client, permitted
credentials\footnote{
  To pass the JWT token back and forth between the server
} and then permitted the two HTTP methods \texttt{POST} and \texttt{GET}.
We also ensured the \texttt{Content-Type, Authorization}, and \texttt{Cookie}
headers are permitted.

Any other methods or headers should abort the request in the preflight.

\subsection{Cookie}\label{subsec:cookie}
JavaScript can access and manipulate \texttt{Cookies}~\cite{he2019malicious}.
We utilized the browser's local cookie storage to attempt secure authentication
between requests.\footnote{
  The cookie storage in a browser is subject to any vulnerabilities that can be
  performed on an SQLite database while having access to the computer where it
  is running.
}
We attempted a couple of strategies listed below to fortify the cookie that
LessPM set in the browser against a malefactor:
\begin{itemize}
  \item \textbf{Strict SameSite}:
  This ensures that the cookie is safeguarded against Cross-Site Request Forgery
  (CSRF) attacks and remains restricted to its original origin domain.
  \item \textbf{Expires}:
  Once the system authenticated the user, it gave the cookie a Time-to-Live
  (TTL) mechanism similar to the JWE, which remained valid for only 15 minutes.
  \item \textbf{Secure}:
  We applied the \texttt{Secure} attribute to ensure that the cookie was only
  accessible through the HTTPS protocol.
  This protocol encrypts the data being sent back and forth between the client
  and the server, attempting to avoid eavesdroppers.
  \item \textbf{HttpOnly}:
  Setting HttpOnly tells the browser to make this cookie inaccessible through
  JavaScript.
  This property is important to mitigate session hijacking.
\end{itemize}

\subsection{Password Encryption}\label{subsec:password-encryption}
A password can be stored and hashed using a \texttt{salt} in a typical
authentication scheme.
When trying to authenticate, the user would provide their UID and password.
These values are collected from the database and checked with the random
salt\footnote{
  Salting is the process of adding a randomly generated string consisting of
  arbitrary characters to the password before creating a hash~\cite{Kharod2015}.
} that was generated when the user registered.

Securing the applicant's passwords made the process more complicated than the
one described above.
Furthermore, since the passwords in a password manager should be a randomly
generated string unknown to the applicant, we cannot hash it and compare the
input from the user.

We decided on the AES-256 (see Section~\ref{subsec:aes}), as a high form of
encryption is paramount to securing the passwords.
AES requires a 256-bit key to encrypt and decrypt the data.
Since the CID generated by WebAuthn is unique and random for each application
and device, this serves as a basis for constructing the key.

We based the key for each password on the following premise:
\begin{enumerate}
  \item
  We take advantage of the fact that each CID is unique in every application
  (see Section~\ref{subsec:webauthn}).
  We therefore used 192-bits of the string that is the CID\@, converting it to
  integers.
  Since every CID is unique depending on the device, we generated padding to
  reach the remaining difference when the CID does not contain 192-bit.
  \item
  We appended each key with a random 128-bit salt of integers and stored these
  bits with the entry in the database.
  \item
  Then we add a 128-bit of pepper collected from the environment variable to
  finish it.
\end{enumerate}
We used this 448-bit as the input for the key-derivation function,
Argon2 (See Section~\ref{subsec:hashing-through-argon2}).
Upon hashing, each key is also subject to a 128-bit salt explicitly generated
for the key.

\subsection{Hashing}\label{subsec:hashing}
\todo[inline]{
When we searched for a good key-derivation function, we first came across
  Password-Based Key Derivation Function 2 (PBKDF2).
We saw PBKDF2 as a good solution for the project, but after researching the
topic further, we ended up with Argon2.

Argon2 is regarded by some to be more secure than PBKDF2 due to its modern
design considerations, including protection against side-channel attacks and
memory-hardness, which make it more resistant to brute-forced and rainbow
table attacks.
Argon2 also offers customization of parameters for flexibility in adapting to
changing hardware technologies and security requirements.
PBKDF2 offers to set an amount of iterations to construct the hash and which
pseudorandom function to use.
}

Since LessPM is required to run in as safe of an environment as possible,
Argon2's configuration option is an excellent solution.

\subsubsection{Configuring Argon2}
Argon2's hashing output is dependable on
configurations~\cite{argon2specs}.\footnote{
  Dependable in this context refers to each configuration that can
  possibly be constructed.
  An instance of Argon2 with 256 Megabytes of \texttt{memory} will not return
  the same hash as 255 Megabytes.
  The same is true for the amount of \texttt{iterations} and
  \texttt{parallelism}.
}
Given that we emphasized security, we opted for the \texttt{Argon2id}
configuration, which gave us equal protection against side-channel and
brute-force attacks.

\todo[inline]{
  A part of Argon2's customizability is the offer to set the option for a
  required amount of memory to do the hashing.
}
These options would force a malefactor to use a specific amount of memory for
each attempt to construct the hash.
The only way a malefactor can get passed this requirement is to purchase more
memory.\footnote{
  As a side note, the increase in memory usage will scale as technology evolves
  and more memory becomes common.
}

For LessPM, we used 128 Megabytes of memory to construct the hash.\footnote{
  According to~\cite{argon2specs}, the reference implementation using Argon2d
  with 4Gb of memory and 8-degree parallelism, the hashing process should take
  0.5 seconds on a CPU with 2Gz. However, we were unsuccessful in seeing
  anywhere close to similar results.
}
We went for the default suggestion of two iterations to complement the memory.
To finalize the configuration, we added 8 degrees of parallelism since the
system where we developed it consists of 8 cores.~\footnote{
  The degree of parallelism is affected by how many cores a CPU
  contains~\cite{argon2specs}.
}

\todo[inline]{The following section}

\subsection{Security Analysis}\label{subsec:security-analysis}
The security of a passwordless password manager is of paramount importance to
protect sensitive user infromation and prevent unauthorized access.
In this subsection, we conduct a comprehensive security analysis, identifying
  potential attack vectors and outlining the defensive measures that have
  been implemented to mitigate tehse risks.

The defensive measures include industry-standard security protocols such as
  HTTPS, JWT/JWE, Argon2 hashing, salt and pepper techniques, CORS
  configuration, and secure cookie management.
Additionally, different JWTs are used for registration, authentication, and
  password creation/retrieval to minimize the risk of exposure.
These defensive measures collectively aim to provide a robust and secure
  architecture for LessPM, safeguarding against various threats and ensuring
  the confidentiality, integrity, and availability of user data.
\begin{enumerate}[label=$\blacktriangleright$]
  \item \textbf{HTTP}

  \textbf{Weakness:}
  Any server running HTTP (Hypertext Transfer Protocol) is passing data
  between a server and a connecting client.
  With HTTP, this data is unencrypted.
  A malefactor can eavesdrop this data, performing a man-in-the-middle attack,
  replacing information, reading tokens, perform header inections, Etc. with 
  this unencrypted traffic.

  \textbf{Defensive Mechanism:}
  Upgrading the HTTP connection to HTTPS ensures that communication between the
  client and server is encrypted through Transport Layer Security (TLS).
  It provides confidentiality and integrity of data transmitted over the
  network, making it more difficult for attackers to intercept or tamper with
  sensitive information.

  \item \textbf{JSON Web Tokens (JWT)}

  \textbf{Weakness:}
  JWTs are signed with a secret specific to the server.
  However, if this secret is discovered or leaked, a malefactor could use this
  information to sign their own tokens and provide their own information.
  Since JWTs are signed and not encrypted, any malefactor who receives access
  to a token can decode the Base64 format and read the tokens information in
  clear text.

  \textbf{Defensive Mechanism:}
  Properly signining the tokens with a strong algorithm prevents information
  from being tampered with.
  An attacker can still read the information in clear-text but can not forge
  their own tokens.
  LessPM uses RSASSA-PSS with Sha-512 to sign and encode tokens, decoding
  them with the keypair when received.
  This ensures integrity and authenticity of the JWT, as the signature is
  verified by the server.
  RSASSA-PSS is a robust and secure signature algorithm that provides protection
  against various attacks, such as collision attacks and length extension
  attacks.
  To further impose a level of security, LessPM takes advantage of different
  tokens for registration, authentication, and authorization.

  \item \textbf{JSON Web Encryption (JWE}

  \textbf{Weakness:}
  JSON Web Encryption are JWTs signed and encrypted using an algorithm for
  signing and encryption.
  However, if this encryption algorithm is not potent or weak (such as a
  Caesar Cipher), an attacker can break the encryption or even encrypt their own
  tokens.
  JWEs are also subject to scenarios such as lack of key rotation, insecure
  length of the encryption key and improper implementation (See
  Section~\ref{sec:futurework}).

  \textbf{Defensive Mechanism:}
  To combat these weaknesses, LessPM implements encryption partly inspired by
  the JWE standard.
  LessPM encrypts the token using AES-256, a strong encryption scheme capable
  of $2^{256}$ different keys.
  This ensures that the data is confidentially kept and protected from
  unauthorized access.
  While inspired, it is important to note that JWEs have not been properly
  implemented in LessPM (See Section~\ref{sec:futurework}).

  \item \textbf{Hashing Passwords \& Keys}

  \textbf{Weakness:}
  Hashing is a one-way process to convert information (such as passwords)
  into a fixed-size string of characters, typically a fixed-length hash value.
  However, the hashing process is susceptible to a malefactor using
  precomputed hash-values for large number of possible passwords, stored in a
  lookup table called a Rainbow table attack.
  Hashes are also vulnerable to dictionary attacks, and not being properly
  implemented without using a salt and preferable a pepper, or brute-force
  attacks.

  \textbf{Defensive Mechanism}
  LessPM takes advantage of the latest within public hashing functionality,
  through Argon2.
  Argon2 constructs a memory-hard and computationally expensive hash that
  provides protection against brute-force, rainbow table, dictionary and
  side-channel attacks by requiring a significant amount of computational
  resources and time to compute the hash.
  Finally, LessPM takes advantage of using both a random salt unique for each
  password and a pepper stored in the source code of 128-bit for both.
  This creates a higher threshold for a malefactor by requiring access to the
  database and the source code to be able to quickly compute a hash.
  The salt and pepper are added to the key of the AES-256

  \item \textbf{Password Authentication}

  \textbf{Weakness}
  A password is something a secret that a user knows.
  Secrets belonging to an individual will always be potentially accessible to
  a malefactor.
  There are multiple vectors that can be used, such as a user being careless
  and writing the password down on a piece of paper or storing the secret in an
  unencrypted file.
  Further, a password is required to be stored somewhere for a server to
  authenticate the user, preferably with an email for recovery.
  This creates two new vectors for a malefactor to exploit, the persisted
  storage on the server or the recovery process.

  \textbf{Defensive Mechanism}
  LessPM takes advantage of WebAuthn to avoid the usage of passwords to
  authenticate the user.
  The user is required to have two vectors for authentication; their
  authenticator device and the biometrics, which is the user itself.
  It is, however, important to mention that any authenticator device is
  vulnerable to being compromised.
  Should a device be compromised or infected with malicious software (Malware),
  it could be used to intercept or manipulate authentication requests.

  \item \textbf{Storing passwords in plaintext}

  \textbf{Weakness:}
  Storing a password can be done in plaintext in a database protected by a
  password.
  However, should an attacker get access to the database, all passwords and
  their plaintext are compromised.
  Storing the passwords in plaintext makes the information accessible to any
  individual with access to the database, which causes a privacy risk.

  \textbf{Defensive Mechanism:}
  LessPM encrypt all passwords with one of the strongest public symmetric
  encryption processes, AES-256.
  Each key for each password is unique, and consists of the 192-bit CID from
  WebAuthn generated when the user registered (See
  Section~\ref{subsec:password-encryption}).
  This CID is unique to each service the user is registered with.
  Along the 192-bit CID, 128-bit is used from a salt and 128-bit used from a
  pepper.
  The salt is stored with the entry in the database.
  Using both salt and pepper enhances the security of the derived key.

  \item \textbf{
    Cross-Site Scripting,
    Cross-Site Request Forgery
    \& Unauthorized Access}

  \textbf{Weakness:}
  Any HTTP server improperly configured is subject Cross-Site Scripting (XSS)
  and CSRF\@.
  XSS is a process where a malefactor could inject a malicious script from a
  different origin, and execute said script in the context of a user's browser.
  This could lead to a situation where a malefactor could get access tokens
  or cookies, or interact with the webpage a user is viewing through the
  script.
  CSRF is a type of security vulnerability where a malefactor tricks a user
  to unknowingly make unauthorized requests on a trusted website, such as a
  bank or similar.
  This can lead to actions being performed on a user's behalf unintentionally,
  without the user's consent, leading to unauthorized access.

  \textbf{Defensive Mechanism:}
  CORS is a measure that enforces a strict policy for which domains and services
  are permitted to access certain resources on the server.
  LessPM takes advantage of CORS by permitting the server itself and the
  client associated domain to access resources on the server.
  This is a preventive measure put in place to allow the communication
  between the client and the serve, even though they are running on different
  ports, and potentially different domains.
  This ensures that only authorized clients can access the server's resources,
  preventing unauthorized cross-origin requests and protecting against
  cross-site scripting (XSS) and cross-site request forgery (CSRF) attacks.

  \item \textbf{Cookies}
  \textbf{Weakness:}
  A cookie is vulnerable to many aspects of security.
  Improperly stored, a cookie can be accessed through JavaScript, sent to
  different domains, get hijacked, excessive expire time, or even poisoned.\footnote{
    Cookie Poisoning is the process where an attacker modifies the content of
    the cookie to inject malicious data or bypass security controls.
  }

  \textbf{Defensive Mechanism:}
  LessPM only uses one cookie, which contains the encrypted JWT\@.
  The cookie is protected through built-in browser-features such as
  restricted to the same origin that the cookie came from, and cannot be sent
  anywhere else, further preventing XSS of sensitive information\@.
  The cookie is expired after 15 minutes, which is extensive amount of time
  for a user to have authorized access to their passwords, before the need to
  reauthenticate their identity.
  Upon creation, LessPM makes sure that the cookie becomes set to secure.
  This prevents the cookie from being sent over an insecure HTTP connection,
  limiting it to HTTPS\@.
  Finally, the cookie is HttpOnly, so that the cookie can't be access through
  JavaScript, reducing the attack vector of SS further.

\end{enumerate}

%\todo[inline]{The preceeding section.}

    \section{Conclusion}\label{sec:conclusion}

    \section{Future work}\label{sec:futurework}
    Through implementing LessPM, we aimed to create a barebone implementation that could serve as a reliable Minimal
    Viable Product (MVP). However, we recognize that more work is needed to further enhance and compliment the product.
    The related topics to further improve LessPM are listed below and briefly discussed as a way to highlight
    potential drawbacks of the current version.

    \begin{enumerate}[label=$\blacktriangleright$]
        \item \textbf{Authorization Headers}
        \item \textbf{AES Key Encryption for Password}
        \item \textbf{Hardcoded AES for JWT}
        \item \textbf{Encrypted Passkey}
        \item \textbf{Attaching connecting IP to JWT}
        \iffalse
        \item \textbf{Credential ID} \newline In its current implementation, LessPM supports a single registered
        authenticator per username to maintain a focused security approach.
        During registration, the server checks the database for an existing username similar to the incoming one
        and aborts the registration ceremony if a match is found.
        While WebAuthn permits users to have multiple authenticators, limiting this feature in the initial iteration of
        LessPM helps ensure a more controlled security environment.
        As the system evolves, considering the addition of support for multiple authenticators can be weighed against
        potential security risks and benefits.
        \fi
    \end{enumerate}
    \printbibliography
\end{document}
